\chapter*{INTRODUCCIÓN}
\addcontentsline{toc}{chapter}{INTRODUCCIÓN}
Las infraestructuras de los proveedores de servicios de internet están compuestas por múltiples dispositivos y sistemas encargados de proporcionar servicios de conectividad y gestionar el tráfico de los usuarios. Estos componentes generan continuamente registros asociados con su operación, autenticación, conectividad, configuración y eventos de seguridad bajo el estandar del protocolo Syslog. Estos registros son fundamentales para la administración de la red, ya que permiten a los administradores monitorear el estado de los dispositivos, identificar problemas, realizar auditorías de seguridad y cumplir con regulaciones y políticas internas.

La diversidad y el volumen de estos registros dificultan su análisis centralizado y la identificación oportuna de eventos que, de manera individual o correlacionada, pueden representar comportamientos anómalos o incidentes de seguridad. Para solucionar este problema se aborda una plataforma de gestión de información y eventos de seguridad (SIEM), la cual permite centralizar y analizar los registros provenientes de todas las fuentes disponibles, automatizando la identificación de patrones y eventos importantes. En el desarrollo del proyecto se plantea una arquitectura mediante las herramientas Syslog-ng y Wazuh, la primera funciona como componente de recepción, procesamiento y normalización de los registros generados por los dispositivos de red, mientras que la segunda se utiliza como plataforma para el análisis de eventos, generación de alertas, correlación y visualización de información de seguridad.

El presente documento tiene como propósito registrar el proceso de diseño, implementación y validación de la plataforma desarrollada, así como servir como referencia técnica para su operación, mantenimiento y futura ampliación. Para ello, se presentan los fundamentos de la solución, la arquitectura implementada, los procedimientos de configuración, los mecanismos de procesamiento y detección de eventos, los resultados obtenidos y las consideraciones para su continuidad.

Este proyecto fue desarrollado en el marco de la práctica profesional realizada en EMCALI EICE ESP, con el propósito de fortalecer las habilidades y complementar la formación académica adquirida a lo largo del programa de pregrado en ingeniería electronica en la Universidad del Valle.